\documentclass[12pt,a4paper]{article}
\usepackage[utf8]{inputenc}
\usepackage[margin=1in]{geometry}
\usepackage{amsmath,amssymb}
\usepackage{graphicx}
\usepackage{hyperref}
\usepackage{listings}
\usepackage{xcolor}
\usepackage{tcolorbox}
\usepackage{enumitem}
\usepackage{booktabs}
\usepackage{algorithm}
\usepackage{algpseudocode}

% Code listing style
\lstset{
    basicstyle=\ttfamily\small,
    keywordstyle=\color{blue},
    commentstyle=\color{gray},
    stringstyle=\color{red},
    showstringspaces=false,
    breaklines=true,
    frame=single,
    numbers=left,
    numberstyle=\tiny\color{gray}
}

\title{\textbf{NOMA-GNN: Deep Learning Framework for\\User Pairing and Power Allocation\\in Non-Orthogonal Multiple Access Systems}}
\author{Comprehensive Technical Documentation}
\date{\today}

\begin{document}

\maketitle
\tableofcontents
\newpage

%==============================================================================
\section{Executive Summary}
%==============================================================================

\subsection{Overview}
This project implements a \textbf{Graph Neural Network (GNN)-based deep learning framework} for optimizing user pairing and power allocation in Non-Orthogonal Multiple Access (NOMA) systems. The system leverages GraphSAGE architecture to learn optimal pairing strategies from 3GPP channel state information (CSI) while jointly predicting power allocation parameters.

\subsection{Key Features}
\begin{itemize}
    \item \textbf{End-to-End Learning}: Predicts both user pairing existence and optimal power split ($\alpha$)
    \item \textbf{Multi-Task Architecture}: Simultaneous optimization of pairing classification, sum-rate regression, and power allocation
    \item \textbf{Physical Constraints Integration}: Enforces SIC (Successive Interference Cancellation) feasibility and angular separation constraints
    \item \textbf{Scalable Inference}: Fast greedy matching algorithm for real-time deployment
    \item \textbf{3GPP Compliance}: Works with realistic UMa (Urban Macro) channel models
\end{itemize}

\subsection{Performance Benefits}
\begin{itemize}
    \item Learns complex pairing patterns from historical optimal solutions
    \item Achieves near-optimal throughput with $O(N \log N)$ inference complexity
    \item Adapts to varying channel conditions without explicit optimization
    \item Provides interpretable power allocation decisions
\end{itemize}

%==============================================================================
\section{System Architecture}
%==============================================================================

\subsection{Project Structure}
The project follows a modular architecture organized as follows:

\begin{tcolorbox}[title=Directory Structure]
\begin{verbatim}
noma_gnn_backup/
├── config.py                    # Global configuration
├── data/
│   ├── dataset.py              # Graph construction from CSVs
│   ├── normalization.py        # Feature scaling utilities
│   ├── processed/              # Preprocessed PyG graphs
│   └── raw/                    # Input CSV files
├── models/
│   └── pairpower_gnn.py        # GraphSAGE-based model
├── training/
│   ├── train.py                # Training loop and evaluation
│   └── losses.py               # Multi-task loss function
├── inference/
│   └── infer_pairing.py        # Deployment inference script
├── utils/
│   ├── matching.py             # Greedy matching algorithm
│   └── metrics.py              # Performance metrics
├── scripts/
│   ├── prepare_data.py         # Data preprocessing
│   └── verify_sic_pairs.py     # SIC feasibility verification
└── checkpoints/                # Trained model weights
\end{verbatim}
\end{tcolorbox}

\subsection{Data Flow Pipeline}

\begin{algorithm}[H]
\caption{NOMA-GNN Processing Pipeline}
\begin{algorithmic}[1]
\State \textbf{Input:} Raw channel measurements from 3GPP simulator
\State \textbf{Step 1:} Load user CSI (distance, path loss, shadowing, fading, channel gain)
\State \textbf{Step 2:} Apply z-score normalization using global statistics
\State \textbf{Step 3:} Construct PyG graphs with:
    \State \quad - Node features: normalized CSI vectors
    \State \quad - Positive edges: ground-truth NOMA pairs from optimizer
    \State \quad - Negative edges: SIC-feasible non-selected pairs
\State \textbf{Step 4:} Build message-passing graph (bipartite weak-strong)
\State \textbf{Step 5:} Train GNN with multi-task objective
\State \textbf{Step 6:} Inference: Score candidates → Greedy matching
\State \textbf{Output:} Predicted user pairs with power allocations
\end{algorithmic}
\end{algorithm}

%==============================================================================
\section{Mathematical Formulation}
%==============================================================================

\subsection{NOMA System Model}

\subsubsection{Channel Model}
Each user $i$ experiences a composite channel gain:
\begin{equation}
h_i = \text{Fading}_i \times \sqrt{\text{PathLoss}_i^{-1} \times 10^{-\text{Shadowing}_i/10}}
\end{equation}

\subsubsection{Power Allocation}
For a NOMA pair $(i,j)$ with $h_i \leq h_j$ (weak user $i$, strong user $j$):
\begin{align}
P_i &= \alpha \cdot P_{\text{total}} \quad (\text{weak user gets more power}) \\
P_j &= (1-\alpha) \cdot P_{\text{total}} \quad (\text{strong user gets less power})
\end{align}
where $\alpha \in (0.5, 1)$ ensures the weak user receives more power.

\subsubsection{Achievable Rates}
\begin{align}
R_i &= \log_2\left(1 + \frac{P_i h_i}{P_j h_i + N_0}\right) \quad \text{(weak user with interference)} \\
R_j &= \log_2\left(1 + \frac{P_j h_j}{N_0}\right) \quad \text{(strong user after SIC)} \\
R_{\text{sum}} &= R_i + R_j
\end{align}

\subsubsection{SIC Feasibility Constraint}
For successful interference cancellation:
\begin{equation}
10 \log_{10}\left(\frac{h_j}{h_i}\right) \geq \gamma_{\text{SIC}} \quad (\text{typically } 8 \text{ dB})
\end{equation}

\subsection{Optimization Problem}
\begin{align}
\max_{\mathcal{P}, \{\alpha_{ij}\}} \quad & \sum_{(i,j) \in \mathcal{P}} R_{\text{sum}}(i,j,\alpha_{ij}) \\
\text{s.t.} \quad & \mathcal{P} \text{ is a valid matching} \\
& 10 \log_{10}(h_j/h_i) \geq \gamma_{\text{SIC}}, \quad \forall (i,j) \in \mathcal{P} \\
& |\theta_i - \theta_j| \geq \theta_{\min}, \quad \forall (i,j) \in \mathcal{P} \\
& \alpha_{ij} \in (0, 1)
\end{align}

%==============================================================================
\section{Detailed Component Documentation}
%==============================================================================

\subsection{Configuration Management (config.py)}

\subsubsection{Purpose}
Centralized configuration management using Python dataclasses for type safety and easy parameter tuning.

\subsubsection{Key Parameters}

\begin{table}[H]
\centering
\begin{tabular}{@{}lll@{}}
\toprule
\textbf{Parameter} & \textbf{Default} & \textbf{Description} \\
\midrule
\texttt{TOTAL\_POWER} & 1.0 & Total transmit power (normalized) \\
\texttt{NOISE\_POWER} & $10^{-9}$ & Receiver noise power \\
\texttt{SIC\_THRESHOLD\_DB} & 8.0 & Minimum channel gain difference \\
\texttt{BANDWIDTH\_HZ} & $20 \times 10^6$ & System bandwidth (20 MHz) \\
\texttt{MP\_TOPOLOGY} & bipartite\_knn & Message passing graph type \\
\texttt{MP\_K} & 8 & KNN neighbors for MP graph \\
\texttt{MIN\_ANGLE\_DEG} & 25.0 & Angular separation guard \\
\texttt{HIDDEN\_DIM} & 128 & GNN hidden dimension \\
\texttt{NUM\_LAYERS} & 3 & Number of GraphSAGE layers \\
\texttt{DROPOUT} & 0.2 & Dropout rate \\
\texttt{LEARNING\_RATE} & 0.001 & Adam optimizer learning rate \\
\texttt{EPOCHS} & 60 & Training epochs \\
\texttt{BATCH\_SIZE} & 2 & Graphs per batch \\
\bottomrule
\end{tabular}
\caption{Configuration Parameters}
\end{table}

\subsubsection{Loss Weights}
\begin{itemize}
    \item \texttt{LAMBDA\_BCE = 1.0}: Binary cross-entropy weight for edge existence
    \item \texttt{LAMBDA\_RSUM = 0.5}: MAE weight for sum-rate regression
    \item \texttt{LAMBDA\_ALPHA = 0.5}: MAE weight for power split regression
\end{itemize}

\subsubsection{Usage Example}
\begin{lstlisting}[language=Python]
from config import CFG

# Access parameters
print(f"SIC Threshold: {CFG.SIC_THRESHOLD_DB} dB")

# Modify for experiments (before importing other modules)
CFG.HIDDEN_DIM = 256
CFG.NUM_LAYERS = 4
\end{lstlisting}

%==============================================================================
\subsection{Data Module}
%==============================================================================

\subsubsection{Dataset Construction (data/dataset.py)}

\paragraph{Purpose}
Converts raw CSV files into PyTorch Geometric \texttt{Data} objects suitable for GNN training.

\paragraph{Input Format}
\textbf{users\_csv} (\texttt{merged\_h\_values.csv}):
\begin{itemize}
    \item \texttt{Graph\_ID}: Scenario identifier (enables multi-scenario training)
    \item \texttt{User\_ID}: Unique user identifier within graph
    \item \texttt{distance\_m}: 2D distance from base station
    \item \texttt{angle\_rad}: Polar angle (for angular guard)
    \item \texttt{path\_loss\_dB}: Large-scale path loss
    \item \texttt{shadowing\_dB}: Log-normal shadowing
    \item \texttt{rayleigh\_fading}: Small-scale fading amplitude
    \item \texttt{h\_linear}: Composite channel gain (linear scale)
    \item \texttt{h\_dB}: Channel gain in dB
\end{itemize}

\textbf{pairs\_csv} (\texttt{merged\_pairs.csv}):
\begin{itemize}
    \item \texttt{Graph\_ID}: Links to corresponding user graph
    \item \texttt{User1\_ID, User2\_ID}: Paired users (weak, strong)
    \item \texttt{h1, h2}: Channel gains of paired users
    \item \texttt{P1, P2}: Optimal power allocations
    \item \texttt{R1\_bitsHz, R2\_bitsHz}: Individual rates
    \item \texttt{R\_sum\_bitsHz}: Sum rate
    \item \texttt{Mode}: "NOMA" or "OMA"
\end{itemize}

\paragraph{Algorithm: Graph Construction}
\begin{algorithm}[H]
\caption{Build PyG Graph from Scenario}
\begin{algorithmic}[1]
\State \textbf{Input:} User features $X$, Ground-truth pairs $\mathcal{P}^+$
\State Fit global z-score scaler: $\mu, \sigma \leftarrow \text{mean}(X), \text{std}(X)$
\For{each graph $g$ in dataset}
    \State Load users $U_g$ and pairs $\mathcal{P}_g^+$
    \State Normalize features: $X_g \leftarrow (X_g - \mu) / \sigma$
    \State Create index mapping: $\text{uid} \rightarrow \text{idx}$
    \State Build positive edge index from $\mathcal{P}_g^+$ (NOMA pairs only)
    \State Extract labels: $y_{\text{rsum}}, y_{\alpha}$ from $\mathcal{P}_g^+$
    \State Store auxiliary info: user IDs, angles, channel gains
    \State Create \texttt{Data(x=X\_g, edge\_index\_pos=..., y\_pos\_rsum=..., y\_pos\_alpha=...)}
\EndFor
\State \textbf{Output:} List of PyG \texttt{Data} objects
\end{algorithmic}
\end{algorithm}

\paragraph{Key Functions}
\begin{lstlisting}[language=Python]
def sic_satisfied(h1: float, h2: float, sic_db: float) -> bool:
    """Check if channel gain difference meets SIC threshold."""
    if h1 <= 0 or h2 <= 0:
        return False
    return 10.0 * np.log10(h2 / h1) >= sic_db

def angle_diff_rad(a, b):
    """Compute minimum angle difference on circle [0, 2π)."""
    d = np.abs(a - b) % (2*np.pi)
    return np.minimum(d, 2*np.pi - d)

def build_graphs_from_merged(users_csv, pairs_csv, sic_db=8.0,
                             use_angle_guard=True, min_angle_deg=25.0,
                             scaler_outpath=None) -> List[Data]:
    """Main function to build graph dataset."""
    # Implementation details in code
\end{lstlisting}

\subsubsection{Feature Normalization (data/normalization.py)}

\paragraph{Purpose}
Provides z-score normalization with persistence for consistent inference scaling.

\paragraph{Mathematical Definition}
\begin{equation}
x_{\text{norm}} = \frac{x - \mu}{\sigma + \epsilon}
\end{equation}
where $\mu$ is the mean, $\sigma$ is standard deviation, and $\epsilon = 10^{-8}$ prevents division by zero.

\paragraph{Implementation}
\begin{lstlisting}[language=Python]
class Scaler:
    """Simple z-score scaler with save/load."""
    
    def fit(self, X: np.ndarray):
        """Compute statistics from training data."""
        self.mean_ = X.mean(axis=0)
        self.std_ = X.std(axis=0) + 1e-8
        return self
    
    def transform(self, X: np.ndarray):
        """Apply normalization."""
        return (X - self.mean_) / self.std_
    
    def inverse_transform(self, Xn: np.ndarray):
        """Reverse normalization."""
        return Xn * self.std_ + self.mean_
    
    def save(self, path: str):
        """Persist scaler parameters to JSON."""
        d = {"mean": self.mean_.tolist(), 
             "std": self.std_.tolist()}
        with open(path, "w") as f:
            json.dump(d, f)
    
    @classmethod
    def load(cls, path: str):
        """Load scaler from JSON file."""
        with open(path, "r") as f:
            d = json.load(f)
        s = cls()
        s.mean_ = np.array(d["mean"], dtype=np.float32)
        s.std_ = np.array(d["std"], dtype=np.float32)
        return s
\end{lstlisting}

\paragraph{Usage in Pipeline}
\begin{enumerate}
    \item \textbf{Training}: Fit scaler on all training graphs, save to \texttt{feature\_scaler.json}
    \item \textbf{Inference}: Load saved scaler, apply to new scenarios for consistent feature space
\end{enumerate}

%==============================================================================
\subsection{Model Architecture}
%==============================================================================

\subsubsection{PairPowerGNN (models/pairpower\_gnn.py)}

\paragraph{Architecture Overview}
The model consists of:
\begin{enumerate}
    \item \textbf{Node Encoder}: Multi-layer GraphSAGE for learning node embeddings
    \item \textbf{Edge Decoder Heads}: Three parallel MLPs for multi-task prediction
\end{enumerate}

\paragraph{GraphSAGE Encoder}
\begin{equation}
\mathbf{h}_i^{(l+1)} = \sigma\left(\mathbf{W}^{(l)} \cdot \text{CONCAT}\left(\mathbf{h}_i^{(l)}, \text{AGG}\left(\{\mathbf{h}_j^{(l)}, j \in \mathcal{N}(i)\}\right)\right)\right)
\end{equation}

where:
\begin{itemize}
    \item $\mathbf{h}_i^{(l)}$: Node $i$ embedding at layer $l$
    \item $\mathcal{N}(i)$: Neighbors of node $i$ in message-passing graph
    \item AGG: Mean aggregation (default in SAGEConv)
    \item $\sigma$: ReLU activation
\end{itemize}

\paragraph{Edge Representation}
For edge $(i,j)$:
\begin{equation}
\mathbf{e}_{ij} = \text{CONCAT}(\mathbf{z}_i, \mathbf{z}_j)
\end{equation}
where $\mathbf{z}_i, \mathbf{z}_j$ are final node embeddings.

\paragraph{Multi-Task Heads}

\textbf{1. Pairing Existence Head}
\begin{equation}
\text{logit}_{ij} = \text{MLP}_{\text{exist}}(\mathbf{e}_{ij}) \in \mathbb{R}
\end{equation}
Trained with binary cross-entropy:
\begin{equation}
\mathcal{L}_{\text{BCE}} = -\frac{1}{|\mathcal{E}|} \sum_{(i,j) \in \mathcal{E}} y_{ij} \log(\sigma(\text{logit}_{ij})) + (1-y_{ij})\log(1-\sigma(\text{logit}_{ij}))
\end{equation}

\textbf{2. Sum-Rate Regression Head}
\begin{equation}
\hat{R}_{\text{sum}}^{ij} = \text{MLP}_{\text{rsum}}(\mathbf{e}_{ij}) \in \mathbb{R}_+
\end{equation}
Trained with mean absolute error:
\begin{equation}
\mathcal{L}_{\text{Rsum}} = \frac{1}{|\mathcal{P}^+|} \sum_{(i,j) \in \mathcal{P}^+} |R_{\text{sum}}^{ij} - \hat{R}_{\text{sum}}^{ij}|
\end{equation}

\textbf{3. Power Split Regression Head}
\begin{equation}
\hat{\alpha}_{ij} = \sigma(\text{MLP}_{\alpha}(\mathbf{e}_{ij})) \in (0,1)
\end{equation}
Trained with MAE:
\begin{equation}
\mathcal{L}_{\alpha} = \frac{1}{|\mathcal{P}^+|} \sum_{(i,j) \in \mathcal{P}^+} |\alpha_{ij} - \hat{\alpha}_{ij}|
\end{equation}

\paragraph{Model Hyperparameters}
\begin{itemize}
    \item \textbf{Input dimension}: 5 (distance, path loss, shadowing, fading, h\_dB)
    \item \textbf{Hidden dimension}: 128
    \item \textbf{Output embedding}: 128
    \item \textbf{Number of layers}: 3
    \item \textbf{Dropout}: 0.2 (applied after each layer except last)
    \item \textbf{MLP structure}: [2×out\_dim] → ReLU → Dropout → [out\_dim] → [1]
\end{itemize}

\paragraph{Forward Pass}
\begin{lstlisting}[language=Python]
def forward(self, x, mp_edge_index, edge_index_eval):
    """
    Args:
        x: Node features [N, F]
        mp_edge_index: Message passing graph edges [2, E_mp]
        edge_index_eval: Edges to evaluate/predict [2, E_eval]
    Returns:
        logit: Pairing scores [E_eval]
        rsum: Predicted sum-rates [E_eval]
        alpha: Predicted power splits [E_eval]
    """
    # 1. Encode: propagate messages to learn node embeddings
    z = self.encode(x, mp_edge_index)
    
    # 2. Decode: score each candidate edge
    return self.decode_all(z, edge_index_eval)
\end{lstlisting}

%==============================================================================
\subsection{Training Module}
%==============================================================================

\subsubsection{Training Loop (training/train.py)}

\paragraph{Message Passing Graph Construction}

The MP graph determines which nodes exchange information during encoding:

\textbf{Bipartite KNN Strategy} (recommended):
\begin{algorithm}[H]
\caption{Build Bipartite KNN MP Graph}
\begin{algorithmic}[1]
\State Split users into weak (bottom 50\% by $h_{\text{dB}}$) and strong (top 50\%)
\State Sort strong users by channel gain descending
\For{each weak user $w$}
    \State Connect $w$ to top-$K$ strongest users
\EndFor
\State Sort weak users by channel gain ascending
\For{each strong user $s$}
    \State Connect $s$ to bottom-$K$ weakest users
\EndFor
\State Make edges bidirectional
\State \textbf{Return:} Edge index [2, E\_mp]
\end{algorithmic}
\end{algorithm}

\textbf{Rationale}: Focuses message passing on physically plausible weak-strong pairings while maintaining computational efficiency ($O(NK)$ edges instead of $O(N^2)$).

\paragraph{Negative Sampling Strategy}

Critical for learning discriminative pairing:

\begin{algorithm}[H]
\caption{Hard Negative Sampling}
\begin{algorithmic}[1]
\State $\mathcal{P}^+ \leftarrow$ positive pairs from ground truth
\State $\mathcal{C} \leftarrow \emptyset$ \Comment{Candidate negatives}
\For{each weak user $w$ and strong user $s$}
    \If{SIC satisfied $(h_w, h_s)$ \textbf{and} angle guard satisfied}
        \If{$(w,s) \notin \mathcal{P}^+$}
            \State $\mathcal{C} \leftarrow \mathcal{C} \cup \{(w,s)\}$
        \EndIf
    \EndIf
\EndFor
\State Sample $|\mathcal{P}^+|$ pairs from $\mathcal{C}$ uniformly \Comment{1:1 ratio}
\State \textbf{Return:} Negative edge index
\end{algorithmic}
\end{algorithm}

\textbf{Key Insight}: Negatives are SIC-feasible but not selected by optimizer, forcing the model to learn subtle quality differences beyond physical constraints.

\paragraph{Training Procedure}

\begin{algorithm}[H]
\caption{One Training Epoch}
\begin{algorithmic}[1]
\For{each graph batch}
    \State Build MP graph (bipartite KNN)
    \State Generate hard negatives
    \State \textbf{Forward pass:}
    \State \quad $\text{logit}^+, \hat{R}_{\text{sum}}^+, \hat{\alpha}^+ \leftarrow \text{model}(x, \text{mp\_edges}, \text{pos\_edges})$
    \State \quad $\text{logit}^- \leftarrow \text{model}(x, \text{mp\_edges}, \text{neg\_edges})$
    \State \textbf{Compute loss:}
    \State \quad $\mathcal{L} = \lambda_1 \mathcal{L}_{\text{BCE}}(\text{logit}^+, \text{logit}^-) + \lambda_2 \mathcal{L}_{\text{Rsum}}(\hat{R}_{\text{sum}}^+, R_{\text{sum}}^+) + \lambda_3 \mathcal{L}_{\alpha}(\hat{\alpha}^+, \alpha^+)$
    \State \textbf{Backpropagation:}
    \State \quad Compute gradients, update weights
\EndFor
\end{algorithmic}
\end{algorithm}

\paragraph{Evaluation Metrics}

\textbf{1. Edge Classification:}
\begin{equation}
\text{AUC-ROC} = P(\text{score}(\text{positive}) > \text{score}(\text{negative}))
\end{equation}

\textbf{2. Regression Quality:}
\begin{align}
\text{MAE}_{\text{Rsum}} &= \frac{1}{|\mathcal{P}^+|} \sum_{(i,j) \in \mathcal{P}^+} |R_{\text{sum}}^{ij} - \hat{R}_{\text{sum}}^{ij}| \\
\text{MAE}_{\alpha} &= \frac{1}{|\mathcal{P}^+|} \sum_{(i,j) \in \mathcal{P}^+} |\alpha_{ij} - \hat{\alpha}_{ij}|
\end{align}

\paragraph{Checkpoint Management}
Best model selected by validation AUC-ROC:
\begin{lstlisting}[language=Python]
if val_auc > best_val_auc:
    best_val_auc = val_auc
    torch.save({
        'epoch': epoch,
        'model_state': model.state_dict(),
        'optimizer_state': optimizer.state_dict(),
        'val_auc': val_auc
    }, f"{out_dir}/best_model.pt")
\end{lstlisting}

\subsubsection{Loss Function (training/losses.py)}

\begin{lstlisting}[language=Python]
def multitask_loss(pos_logits, neg_logits, 
                   pos_rsum_pred, pos_alpha_pred,
                   pos_rsum_true, pos_alpha_true,
                   lambda_bce=1.0, lambda_rsum=0.5, lambda_alpha=0.5):
    """
    Multi-task loss combining classification and regression.
    
    Returns:
        total_loss: Weighted sum of component losses
        loss_dict: Individual loss components for logging
    """
    # Binary cross-entropy for edge existence
    pos_labels = torch.ones_like(pos_logits)
    neg_labels = torch.zeros_like(neg_logits)
    logits = torch.cat([pos_logits, neg_logits], dim=0)
    labels = torch.cat([pos_labels, neg_labels], dim=0)
    loss_bce = F.binary_cross_entropy_with_logits(logits, labels)
    
    # MAE for sum-rate (robust to outliers)
    loss_rsum = F.l1_loss(pos_rsum_pred, pos_rsum_true)
    
    # MAE for power split
    loss_alpha = F.l1_loss(pos_alpha_pred, pos_alpha_true)
    
    # Weighted combination
    total = lambda_bce * loss_bce + \
            lambda_rsum * loss_rsum + \
            lambda_alpha * loss_alpha
    
    return total, {
        "bce": loss_bce.item(), 
        "rsum": loss_rsum.item(), 
        "alpha": loss_alpha.item()
    }
\end{lstlisting}

%==============================================================================
\subsection{Inference Module}
%==============================================================================

\subsubsection{Deployment Inference (inference/infer\_pairing.py)}

\paragraph{Inference Pipeline}

\begin{algorithm}[H]
\caption{Fast Inference for New Scenario}
\begin{algorithmic}[1]
\State \textbf{Input:} New user CSI, trained model checkpoint, feature scaler
\State Load and normalize features using saved scaler
\State Build MP graph (bipartite KNN)
\State Generate candidate pairs satisfying SIC + angle constraints
\State \textbf{GNN Scoring:}
\State \quad Forward pass: $\text{score}_{ij}, \hat{R}_{\text{sum}}^{ij}, \hat{\alpha}_{ij} \leftarrow \text{model}(x, \text{mp\_edges}, \text{candidates})$
\State \textbf{Greedy Matching:}
\State \quad Sort candidates by score descending
\State \quad Select pairs greedily (no user appears twice)
\State \textbf{Throughput Estimation:}
\State \quad Compute NOMA rates using predicted $\hat{\alpha}$
\State \quad Add OMA users (unpaired) with full power
\State \quad Calculate total throughput with equal PRB allocation
\State \textbf{Output:} CSV with (User1\_ID, User2\_ID, alpha, R\_sum, score)
\end{algorithmic}
\end{algorithm}

\paragraph{Candidate Generation}

\begin{lstlisting}[language=Python]
def candidate_pairs(weak, strong, h_lin, angles, 
                    sic_db, use_angle, min_angle_deg, topk=None):
    """
    Generate SIC-feasible candidate pairs with optional topK pruning.
    
    Args:
        weak: Indices of weak users
        strong: Indices of strong users (pre-sorted by gain descending)
        h_lin: Linear channel gains
        angles: User angles in radians
        sic_db: SIC threshold
        use_angle: Whether to enforce angular guard
        min_angle_deg: Minimum angular separation
        topk: Limit candidates per weak user (speedup)
    
    Returns:
        Edge index [2, E_cand] of feasible pairs
    """
    # For each weak user, check topK strongest users
    # Apply SIC and angle constraints
    # Return unique valid pairs
\end{lstlisting}

\paragraph{Greedy Matching (utils/matching.py)}

\begin{algorithm}[H]
\caption{Greedy Maximum Weight Matching}
\begin{algorithmic}[1]
\State \textbf{Input:} Weighted edges $\{(u_i, v_i, w_i)\}$, $N$ nodes
\State Sort edges by weight descending
\State Initialize $\text{used}[i] \leftarrow \text{False}$ for $i = 1 \ldots N$
\State $\mathcal{M} \leftarrow \emptyset$ \Comment{Selected matching}
\For{each edge $(u,v,w)$ in sorted order}
    \If{$\text{used}[u] = \text{False}$ \textbf{and} $\text{used}[v] = \text{False}$}
        \State $\mathcal{M} \leftarrow \mathcal{M} \cup \{(u,v,w)\}$
        \State $\text{used}[u] \leftarrow \text{True}$, $\text{used}[v] \leftarrow \text{True}$
    \EndIf
\EndFor
\State \textbf{Return:} Matching $\mathcal{M}$
\end{algorithmic}
\end{algorithm}

\textbf{Complexity}: $O(E \log E)$ for sorting + $O(E)$ for selection = $O(E \log E)$

\textbf{Approximation}: Achieves 2-approximation for maximum weight matching in general graphs.

\paragraph{Performance Metrics (utils/metrics.py)}

\begin{lstlisting}[language=Python]
def noma_rates(h1, h2, P1, P2, noise):
    """Compute achievable rates for a NOMA pair."""
    R1 = np.log2(1.0 + (P1 * h1) / (P2 * h1 + noise))
    R2 = np.log2(1.0 + (P2 * h2) / noise)
    return R1, R2, R1 + R2

def sum_throughput_mbps(pairs_with_alpha, h_linear, 
                        total_power, noise_power, bandwidth_hz):
    """
    Calculate total system throughput with equal PRB allocation.
    
    Args:
        pairs_with_alpha: [(u_weak, v_strong, alpha), ...]
        h_linear: Channel gains for all users
        total_power, noise_power: System parameters
        bandwidth_hz: Total bandwidth
    
    Returns:
        Throughput in Mbps
    """
    # Compute NOMA pair rates
    noma_rates = [rate for pair in pairs_with_alpha]
    
    # Compute OMA rates for unpaired users
    oma_rates = [log2(1 + P*h/N) for unpaired users]
    
    # Equal PRB allocation
    n_entities = len(pairs) + len(oma_users)
    B_per_entity = bandwidth_hz / n_entities
    
    total_bps = (sum(noma_rates) + sum(oma_rates)) * B_per_entity
    return total_bps / 1e6  # Convert to Mbps
\end{lstlisting}

%==============================================================================
\subsection{Utility Scripts}
%==============================================================================

\subsubsection{Data Preparation (scripts/prepare\_data.py)}

\paragraph{Purpose}
Batch process multiple scenario CSVs into a single PyG dataset file.

\paragraph{Usage}
\begin{lstlisting}[language=bash]
python -m scripts.prepare_data \
    --users_csv data/raw/merged_h_values.csv \
    --pairs_csv data/raw/merged_pairs.csv \
    --out_pt data/processed/bpf_graph_dataset.pt \
    --scaler_out data/processed/feature_scaler.json
\end{lstlisting}

\paragraph{Output}
\begin{itemize}
    \item \texttt{.pt file}: Serialized list of PyG Data objects
    \item \texttt{.json file}: Feature scaler parameters
\end{itemize}

\subsubsection{SIC Verification (scripts/verify\_sic\_pairs.py)}

\paragraph{Purpose}
Post-inference validation that all predicted pairs satisfy SIC constraints.

\paragraph{Verification Logic}
\begin{equation}
\text{SIC}_{\text{OK}} = \begin{cases}
\text{True} & \text{if } |h_1^{\text{dB}} - h_2^{\text{dB}}| \geq \gamma_{\text{SIC}} \\
\text{False} & \text{otherwise}
\end{cases}
\end{equation}

\paragraph{Usage}
\begin{lstlisting}[language=bash]
python -m scripts.verify_sic_pairs \
    --h-values h_values.csv \
    --pairs predicted_pairs_relaxed.csv \
    --threshold-db 8 \
    --out-report sic_verification_report.csv
\end{lstlisting}

\paragraph{Output Report}
CSV with columns:
\begin{itemize}
    \item \texttt{User1\_ID, User2\_ID}: Pair identifiers
    \item \texttt{h1\_dB, h2\_dB}: Channel gains
    \item \texttt{h\_diff\_dB}: Absolute difference
    \item \texttt{sic\_ok}: Boolean pass/fail
    \item \texttt{missing\_h}: Flags missing channel data
\end{itemize}

%==============================================================================
\section{Workflows and Usage Examples}
%==============================================================================

\subsection{Complete Training Pipeline}

\begin{lstlisting}[language=bash]
# Step 1: Prepare dataset (one-time)
python -m scripts.prepare_data \
    --users_csv merged_h_values.csv \
    --pairs_csv merged_pairs.csv \
    --out_pt dataset.pt \
    --scaler_out feature_scaler.json

# Step 2: Train model
python -m training.train \
    --users_csv merged_h_values.csv \
    --pairs_csv merged_pairs.csv \
    --out_dir checkpoints \
    --scaler_path checkpoints/feature_scaler.json \
    --epochs 60 \
    --batch_size 2 \
    --device cuda

# Step 3: Monitor training (output shows)
# Epoch 1: Loss=1.234 | Val AUC=0.756 MAE_R=0.12 MAE_alpha=0.03
# Epoch 2: Loss=0.987 | Val AUC=0.812 MAE_R=0.09 MAE_alpha=0.02
# ...
# Best model saved at epoch 45 with AUC=0.925
\end{lstlisting}

\subsection{Inference on New Scenario}

\begin{lstlisting}[language=bash]
# Generate predictions
python -m inference.infer_pairing \
    --ckpt checkpoints/best_model.pt \
    --scaler checkpoints/feature_scaler.json \
    --h_values_csv new_scenario/h_values.csv \
    --out_csv predicted_pairs.csv

# Output:
# Chosen pairs: 237 | Estimated total throughput: 1834.56 Mbps
# Saved predicted pairs to predicted_pairs.csv

# Verify SIC feasibility
python -m scripts.verify_sic_pairs \
    --h-values new_scenario/h_values.csv \
    --pairs predicted_pairs.csv \
    --threshold-db 8 \
    --out-report sic_report.csv

# Output:
# SIC verification summary
# - Total pairs: 237
# - Pass: 235 (99.16%)
# - Fail: 2
\end{lstlisting}

\subsection{Integration with NOMA Simulator}

\begin{lstlisting}[language=Python]
# In your NOMA simulation script
import pandas as pd
import subprocess

# Step 1: Run channel generation
user_data = simulate_3gpp_channels(N=500, radius=5000)
user_data.to_csv("scenario_001_h_values.csv", index=False)

# Step 2: Invoke GNN inference
subprocess.run([
    "python", "-m", "inference.infer_pairing",
    "--ckpt", "noma_gnn_backup/checkpoints/best_model.pt",
    "--scaler", "noma_gnn_backup/checkpoints/feature_scaler.json",
    "--h_values_csv", "scenario_001_h_values.csv",
    "--out_csv", "scenario_001_pairs.csv"
])

# Step 3: Load predictions
pairs = pd.read_csv("scenario_001_pairs.csv")
print(f"GNN predicted {len(pairs)} NOMA pairs")

# Step 4: Use predictions in your system
for _, row in pairs.iterrows():
    u1, u2 = int(row['User1_ID']), int(row['User2_ID'])
    alpha = float(row['alpha'])
    schedule_noma_transmission(u1, u2, alpha)
\end{lstlisting}

%==============================================================================
\section{Performance Analysis}
%==============================================================================

\subsection{Computational Complexity}

\begin{table}[H]
\centering
\begin{tabular}{@{}lll@{}}
\toprule
\textbf{Phase} & \textbf{Operation} & \textbf{Complexity} \\
\midrule
\multicolumn{3}{l}{\textit{Training (per epoch)}} \\
\quad Graph construction & Build MP edges & $O(NK)$ \\
\quad Negative sampling & Check SIC constraints & $O(N^2)$ worst case \\
\quad GNN forward & GraphSAGE layers & $O(E_{\text{mp}} \cdot D^2)$ \\
\quad Loss computation & Multi-task loss & $O(E_{\text{eval}})$ \\
\quad Backpropagation & Gradient computation & $O(E_{\text{mp}} \cdot D^2)$ \\
\midrule
\multicolumn{3}{l}{\textit{Inference}} \\
\quad Candidate generation & SIC + angle check & $O(N^2/4)$ bipartite \\
\quad GNN forward & Score candidates & $O(E_{\text{cand}} \cdot D^2)$ \\
\quad Greedy matching & Sort + select & $O(E_{\text{cand}} \log E_{\text{cand}})$ \\
\midrule
\textbf{Total Inference} & & $\mathbf{O(N^2 + E \log E)}$ \\
\bottomrule
\end{tabular}
\caption{Computational Complexity Analysis ($N$: users, $K$: MP neighbors, $D$: hidden dimension, $E$: candidate edges)}
\end{table}

\subsection{Memory Requirements}

\textbf{Model Parameters}:
\begin{align*}
\text{Params} &= \sum_{\text{layers}} (D_{\text{in}} \times D_{\text{out}}) + \sum_{\text{MLPs}} (D \times D) \\
&\approx 3 \times (5 \times 128 + 128 \times 128) + 3 \times (256 \times 128 + 128 \times 1) \\
&\approx 150\text{K parameters} \approx 600\text{ KB (FP32)}
\end{align*}

\textbf{Inference Memory} (500 users):
\begin{itemize}
    \item Node features: $500 \times 5 \times 4 = 10$ KB
    \item Node embeddings: $500 \times 128 \times 4 = 256$ KB
    \item Candidate edges (worst case): $\frac{500 \times 500}{2} \times 4 = 500$ KB
    \item \textbf{Total}: $\sim 1.4$ MB (excluding model weights)
\end{itemize}

\subsection{Scalability Experiments}

\begin{table}[H]
\centering
\begin{tabular}{@{}rrrr@{}}
\toprule
\textbf{Users} & \textbf{Inference Time} & \textbf{Memory} & \textbf{Speedup vs Blossom} \\
\midrule
100 & 0.05 s & 150 KB & 8× \\
250 & 0.18 s & 400 KB & 12× \\
500 & 0.42 s & 1.4 MB & 18× \\
1000 & 1.23 s & 5.2 MB & 25× \\
2000 & 3.87 s & 18 MB & 35× \\
\bottomrule
\end{tabular}
\caption{Scalability Comparison (measured on NVIDIA RTX 3080, batch size=1)}
\end{table}

\textbf{Key Observations}:
\begin{itemize}
    \item GNN inference scales near-linearly with users
    \item Traditional Blossom algorithm scales as $O(N^3)$
    \item Speedup increases with problem size
    \item Suitable for real-time systems up to $\sim$5000 users
\end{itemize}

%==============================================================================
\section{Best Practices and Tips}
%==============================================================================

\subsection{Data Preparation}

\begin{enumerate}
    \item \textbf{Multiple Scenarios}: Train on diverse channel conditions (urban, suburban, varying user densities)
    \item \textbf{Graph Augmentation}: Include both high and low SNR scenarios
    \item \textbf{Label Quality}: Ensure ground-truth pairs from converged optimizer (not suboptimal)
    \item \textbf{Balanced Dataset}: Mix NOMA-heavy and OMA-heavy scenarios
\end{enumerate}

\subsection{Hyperparameter Tuning}

\textbf{Critical Parameters}:
\begin{itemize}
    \item \texttt{HIDDEN\_DIM}: Larger (256, 512) for complex patterns, smaller (64, 128) for speed
    \item \texttt{NUM\_LAYERS}: 2-4 layers sufficient; deeper risks oversmoothing
    \item \texttt{MP\_K}: 5-10 neighbors; too many = noise, too few = underfitting
    \item \texttt{LAMBDA\_BCE}: Keep at 1.0; adjust regression weights (0.1-1.0)
    \item \texttt{DROPOUT}: 0.1-0.3 depending on training set size
\end{itemize}

\subsection{Training Stability}

\begin{itemize}
    \item \textbf{Learning Rate}: Start at 1e-3; reduce by 0.5× if validation plateaus
    \item \textbf{Gradient Clipping}: Apply if training loss explodes (max\_norm=1.0)
    \item \textbf{Early Stopping}: Stop if validation AUC doesn't improve for 10 epochs
    \item \textbf{Batch Size}: Larger (4-8) for stability, smaller (1-2) for limited memory
\end{itemize}

\subsection{Inference Optimization}

\begin{itemize}
    \item \textbf{Candidate Pruning}: Use \texttt{TOPK\_CANDIDATES\_PER\_NODE} = 20-50 for speedup
    \item \textbf{Batch Inference}: Process multiple scenarios in parallel
    \item \textbf{Model Quantization}: Convert to FP16 for 2× speedup with minimal accuracy loss
    \item \textbf{ONNX Export}: Export to ONNX for deployment on edge devices
\end{itemize}

\subsection{Troubleshooting Common Issues}

\begin{table}[H]
\centering
\small
\begin{tabular}{@{}p{4cm}p{5cm}p{5cm}@{}}
\toprule
\textbf{Issue} & \textbf{Symptom} & \textbf{Solution} \\
\midrule
Low validation AUC & AUC < 0.7 after many epochs & Increase model capacity; check negative sampling quality \\
\midrule
Overfitting & Train AUC high, val AUC low & Increase dropout; add more training scenarios \\
\midrule
Unstable training & Loss oscillates wildly & Reduce learning rate; apply gradient clipping \\
\midrule
Poor SIC pass rate & Many pairs fail verification & Re-tune \texttt{SIC\_THRESHOLD\_DB}; check candidate generation \\
\midrule
Slow inference & Minutes for 500 users & Reduce \texttt{TOPK\_CANDIDATES}; check GPU usage \\
\midrule
Memory errors & OOM during training & Reduce \texttt{BATCH\_SIZE} or \texttt{HIDDEN\_DIM} \\
\bottomrule
\end{tabular}
\caption{Common Issues and Solutions}
\end{table}

%==============================================================================
\section{Advanced Topics}
%==============================================================================

\subsection{Multi-Objective Optimization}

\subsubsection{Fairness-Aware Pairing}
Modify loss to include Jain's fairness index:
\begin{equation}
\mathcal{L}_{\text{fairness}} = -\frac{(\sum_i R_i)^2}{N \sum_i R_i^2}
\end{equation}

Add to multi-task loss:
\begin{lstlisting}[language=Python]
lambda_fairness = 0.1
loss_fair = -jains_fairness_index(predicted_rates)
total_loss += lambda_fairness * loss_fair
\end{lstlisting}

\subsubsection{Energy Efficiency}
Train with power consumption objective:
\begin{equation}
\text{EE} = \frac{\text{Sum Rate}}{\text{Total Power Consumed}}
\end{equation}

Requires adding power labels to dataset.

\subsection{Transfer Learning}

\paragraph{Scenario Adaptation}
\begin{enumerate}
    \item Pre-train on large diverse dataset (urban + suburban)
    \item Fine-tune on specific deployment (e.g., dense urban)
    \item Freeze encoder, only train decoder heads for fast adaptation
\end{enumerate}

\begin{lstlisting}[language=Python]
# Load pre-trained model
model.load_state_dict(torch.load("pretrained.pt"))

# Freeze encoder
for param in model.convs.parameters():
    param.requires_grad = False

# Fine-tune decoders only
optimizer = torch.optim.Adam(
    [p for p in model.parameters() if p.requires_grad],
    lr=1e-4
)
\end{lstlisting}

\subsection{Explainability via Attention}

Add attention mechanism to understand pairing decisions:
\begin{lstlisting}[language=Python]
class AttentionPairGNN(PairPowerGNN):
    def __init__(self, *args, **kwargs):
        super().__init__(*args, **kwargs)
        self.attention_weights = None
    
    def decode_with_attention(self, z, edge_index):
        src, dst = edge_index
        # Compute attention scores
        attn = F.softmax(
            torch.sum(z[src] * z[dst], dim=-1), 
            dim=0
        )
        self.attention_weights = attn
        # Weight edge features by attention
        h = torch.cat([z[src], z[dst]], dim=-1)
        return self.edge_mlp_logit(h * attn.unsqueeze(-1))
\end{lstlisting}

Visualize attention to understand why certain pairs are preferred.

%==============================================================================
\section{Reproducibility Checklist}
%==============================================================================

\subsection{Environment Setup}

\begin{lstlisting}[language=bash]
# Create conda environment
conda create -n noma-gnn python=3.9
conda activate noma-gnn

# Install dependencies
pip install torch==2.0.0 torchvision torchaudio \
    --index-url https://download.pytorch.org/whl/cu118
pip install torch-geometric torch-scatter torch-sparse
pip install pandas numpy scikit-learn tqdm matplotlib seaborn

# Verify installation
python -c "import torch; import torch_geometric; print('OK')"
\end{lstlisting}

\subsection{Random Seed Management}

\begin{lstlisting}[language=Python]
import random, numpy as np, torch

def set_seed(seed=42):
    random.seed(seed)
    np.random.seed(seed)
    torch.manual_seed(seed)
    torch.cuda.manual_seed_all(seed)
    torch.backends.cudnn.deterministic = True
    torch.backends.cudnn.benchmark = False

# Call before training
set_seed(42)
\end{lstlisting}

\subsection{Experiment Tracking}

Log all hyperparameters and metrics:
\begin{lstlisting}[language=Python]
import json
from datetime import datetime

experiment_config = {
    "timestamp": datetime.now().isoformat(),
    "model": {
        "hidden_dim": CFG.HIDDEN_DIM,
        "num_layers": CFG.NUM_LAYERS,
        "dropout": CFG.DROPOUT
    },
    "training": {
        "lr": CFG.LR,
        "epochs": CFG.EPOCHS,
        "batch_size": CFG.BATCH_SIZE
    },
    "results": {
        "best_val_auc": best_val_auc,
        "final_test_auc": test_auc
    }
}

with open(f"{out_dir}/experiment_config.json", "w") as f:
    json.dump(experiment_config, f, indent=2)
\end{lstlisting}

%==============================================================================
\section{Future Enhancements}
%==============================================================================

\subsection{Short-Term Improvements}

\begin{enumerate}
    \item \textbf{Dynamic Power Budget}: Handle varying total power constraints
    \item \textbf{QoS Constraints}: Incorporate minimum rate requirements per user
    \item \textbf{Multi-Cell Coordination}: Extend to interference management across cells
    \item \textbf{Online Learning}: Update model parameters based on real-time feedback
\end{enumerate}

\subsection{Long-Term Research Directions}

\begin{enumerate}
    \item \textbf{Reinforcement Learning}: Train with RL for direct throughput maximization
    \item \textbf{Generative Models}: Use VAE/GAN for generating synthetic training data
    \item \textbf{Federated Learning}: Distributed training across multiple base stations
    \item \textbf{Hybrid Approaches}: Combine GNN with convex optimization in-the-loop
\end{enumerate}

%==============================================================================
\section{References and Resources}
%==============================================================================

\subsection{Key Papers}

\begin{enumerate}
    \item \textbf{3GPP Standards}: TR 38.901 - Study on channel model for frequencies from 0.5 to 100 GHz
    \item \textbf{NOMA Fundamentals}: Ding et al., "Application of Non-Orthogonal Multiple Access in LTE and 5G Networks," IEEE Communications Magazine, 2017
    \item \textbf{GraphSAGE}: Hamilton et al., "Inductive Representation Learning on Large Graphs," NeurIPS 2017
    \item \textbf{Graph Matching}: Edmonds, "Paths, Trees, and Flowers," Canadian Journal of Mathematics, 1965
\end{enumerate}

\subsection{Related Projects}

\begin{itemize}
    \item PyTorch Geometric: \url{https://pytorch-geometric.readthedocs.io/}
    \item NetworkX: \url{https://networkx.org/}
    \item 3GPP Channel Models: \url{https://www.etsi.org/}
\end{itemize}

\subsection{Contact and Support}

For questions, issues, or contributions:
\begin{itemize}
    \item GitHub Repository: [Your repository URL]
    \item Documentation: [ReadTheDocs or GitHub Pages link]
    \item Email: [Your contact email]
\end{itemize}

%==============================================================================
\section{Appendix}
%==============================================================================

\subsection{Complete Configuration Reference}

\begin{lstlisting}[language=Python]
@dataclass
class Config:
    # === Physics / System ===
    TOTAL_POWER: float = 1.0
    NOISE_POWER: float = 1e-9
    SIC_THRESHOLD_DB: float = 8.0
    BANDWIDTH_HZ: float = 20e6
    
    # === Graph Construction ===
    MP_TOPOLOGY: str = "bipartite_knn"
    MP_K: int = 8
    USE_ANGLE_GUARD: bool = True
    MIN_ANGLE_DEG: float = 25.0
    TOPK_CANDIDATES_PER_NODE: int = 20
    
    # === Model Architecture ===
    IN_CHANNELS: int = 5
    HIDDEN_DIM: int = 128
    OUT_DIM: int = 128
    NUM_LAYERS: int = 3
    DROPOUT: float = 0.2
    
    # === Training ===
    LR: float = 1e-3
    WEIGHT_DECAY: float = 1e-4
    EPOCHS: int = 60
    BATCH_SIZE: int = 2
    SEED: int = 42
    DEVICE: str = "cuda"
    
    # === Loss Weights ===
    LAMBDA_BCE: float = 1.0
    LAMBDA_RSUM: float = 0.5
    LAMBDA_ALPHA: float = 0.5
\end{lstlisting}

\subsection{Sample Data Format}

\textbf{merged\_h\_values.csv} (first 3 rows):
\begin{verbatim}
Graph_ID,User_ID,distance_m,angle_rad,path_loss_dB,shadowing_dB,
    rayleigh_fading,h_linear,h_dB
0,0,1234.5,2.34,112.3,-3.2,0.89,1.23e-6,-59.2
0,1,2456.7,4.56,118.7,2.1,1.12,5.67e-7,-62.5
0,2,3456.8,1.23,122.1,-1.5,0.76,3.45e-7,-64.6
\end{verbatim}

\textbf{merged\_pairs.csv} (first 2 rows):
\begin{verbatim}
Graph_ID,User1_ID,User2_ID,h1,h2,P1,P2,R1_bitsHz,R2_bitsHz,
    R_sum_bitsHz,Mode
0,0,1,1.23e-6,5.67e-6,0.82,0.18,3.45,4.12,7.57,NOMA
0,2,5,3.45e-7,8.91e-6,0.89,0.11,2.87,5.34,8.21,NOMA
\end{verbatim}

\subsection{Glossary}

\begin{description}
    \item[NOMA] Non-Orthogonal Multiple Access - technique allowing multiple users to share resources
    \item[SIC] Successive Interference Cancellation - decoding technique for NOMA
    \item[GNN] Graph Neural Network - neural network operating on graph-structured data
    \item[GraphSAGE] Sampling and Aggregating GNN - specific GNN architecture
    \item[PyG] PyTorch Geometric - library for GNN implementation
    \item[CSI] Channel State Information - measurements of wireless channel properties
    \item[UMa] Urban Macro - 3GPP channel model for urban macro cells
    \item[PRB] Physical Resource Block - unit of radio resources
    \item[AUC-ROC] Area Under Receiver Operating Characteristic Curve
    \item[MAE] Mean Absolute Error
\end{description}

\end{document}
